\documentclass[10pt,letterpaper,final]{report}
\usepackage[margin=.75in]{geometry}
\usepackage[utf8]{inputenc}
\usepackage{amsmath}
\usepackage{amsfonts}
\usepackage{amssymb}
\usepackage{amsthm}
\renewcommand\qedsymbol{$\blacksquare$}
\usepackage{enumerate}
\usepackage{enumitem}
\usepackage{hyperref}
\usepackage{pdfpages}
\usepackage{graphics}
\usepackage{graphicx}
\usepackage{tikz}
\usepackage{tikz-qtree}
\usepackage{forest}
\usepackage{float}
\usetikzlibrary{automata,arrows}
\usetikzlibrary{shapes.multipart, positioning}
\usepackage{mdframed}
\usepackage[
  backend=biber,
  style=numeric-comp,
  sorting=none
]{biblatex}
\addbibresource{references.bib}

\usepackage{minted}
\usemintedstyle{algol_nu}
\usepackage{xltabular}
\usepackage{pdflscape}
\usepackage{tcolorbox}
\usepackage{enumitem}


% Based on template from COSC-417 at Towson University, originally authored by Marius Zimand

\begin{document}
\begin{center}
  \fbox{
    \vbox{
      \begin{flushleft}
        Jonathan Llovet \\  % authors' names
        EN.605.202.81 \\  %class
        2026-08-11 \\  % date
        Module 10 - Homework 10 \\ % ID
      \end{flushleft}
      \center{\Large{\textbf{Hashing}}}
      %\end{mdframed}
    } % end vbox
  } % end fbox
  \vline
\end{center}

\medskip

For the following problem sets, you will perform hashing using different collision mechanisms.  The hashtable has a size of 13 and will use the following hashing algorithm that returns an integer:

\begin{minted}[breaklines,breakanywhere=true,fontsize={\fontsize{8pt}{9pt}\selectfont}]{python}
def get_hashed_value(key) -> int:
    return key * 2 + 3
\end{minted}

For each of the problem sets, you will keep track of the number of collisions and show your work for each step after inserting the following values:

\texttt{5, 4, 25, 8, 10, 34, 18, 51, 17, 21}

\medskip
\noindent\textbf{Problem 1:}
\medskip

Linear Probing - For this example, use ``simple'' linear probing in order to insert the values into the hash table.

\medskip
\noindent\textbf{Answer:}
\medskip

A hash table with linear probing handles a collision by starting at the key's mapped bucket, and then linearly searches subsequent buckets until an empty bucket is found. \parencite{zybooks-ds2026}

\begin{verbatim}
Inserting 5 into ['_', '_', '_', '_', '_', '_', '_', '_', '_', '_', '_', '_', '_']
get_hashed_value(5) -> 13 -> x mod 13 -> 0
Attempting to insert 5 at index 0
Inserted 5 at index 0
0 collisions while inserting 5
Current value of hash table: [5, '_', '_', '_', '_', '_', '_', '_', '_', '_', '_', '_', '_']

Inserting 4 into [5, '_', '_', '_', '_', '_', '_', '_', '_', '_', '_', '_', '_']
get_hashed_value(4) -> 11 -> x mod 13 -> 11
Attempting to insert 4 at index 11
Inserted 4 at index 11
0 collisions while inserting 4
Current value of hash table: [5, '_', '_', '_', '_', '_', '_', '_', '_', '_', '_', 4, '_']

Inserting 25 into [5, '_', '_', '_', '_', '_', '_', '_', '_', '_', '_', 4, '_']
get_hashed_value(25) -> 53 -> x mod 13 -> 1
Attempting to insert 25 at index 1
Inserted 25 at index 1
0 collisions while inserting 25
Current value of hash table: [5, 25, '_', '_', '_', '_', '_', '_', '_', '_', '_', 4, '_']

Inserting 8 into [5, 25, '_', '_', '_', '_', '_', '_', '_', '_', '_', 4, '_']
get_hashed_value(8) -> 19 -> x mod 13 -> 6
Attempting to insert 8 at index 6
Inserted 8 at index 6
0 collisions while inserting 8
Current value of hash table: [5, 25, '_', '_', '_', '_', 8, '_', '_', '_', '_', 4, '_']

Inserting 10 into [5, 25, '_', '_', '_', '_', 8, '_', '_', '_', '_', 4, '_']
get_hashed_value(10) -> 23 -> x mod 13 -> 10
Attempting to insert 10 at index 10
Inserted 10 at index 10
0 collisions while inserting 10
Current value of hash table: [5, 25, '_', '_', '_', '_', 8, '_', '_', '_', 10, 4, '_']

Inserting 34 into [5, 25, '_', '_', '_', '_', 8, '_', '_', '_', 10, 4, '_']
get_hashed_value(34) -> 71 -> x mod 13 -> 6
Attempting to insert 34 at index 6
Collision at 6, which contains the value: 8
Attempting to insert 34 at index 7
Inserted 34 at index 7
1 collisions while inserting 34
Current value of hash table: [5, 25, '_', '_', '_', '_', 8, 34, '_', '_', 10, 4, '_']

Inserting 18 into [5, 25, '_', '_', '_', '_', 8, 34, '_', '_', 10, 4, '_']
get_hashed_value(18) -> 39 -> x mod 13 -> 0
Attempting to insert 18 at index 0
Collision at 0, which contains the value: 5
Attempting to insert 18 at index 1
Collision at 1, which contains the value: 25
Attempting to insert 18 at index 2
Inserted 18 at index 2
2 collisions while inserting 18
Current value of hash table: [5, 25, 18, '_', '_', '_', 8, 34, '_', '_', 10, 4, '_']

Inserting 51 into [5, 25, 18, '_', '_', '_', 8, 34, '_', '_', 10, 4, '_']
get_hashed_value(51) -> 105 -> x mod 13 -> 1
Attempting to insert 51 at index 1
Collision at 1, which contains the value: 25
Attempting to insert 51 at index 2
Collision at 2, which contains the value: 18
Attempting to insert 51 at index 3
Inserted 51 at index 3
2 collisions while inserting 51
Current value of hash table: [5, 25, 18, 51, '_', '_', 8, 34, '_', '_', 10, 4, '_']

Inserting 17 into [5, 25, 18, 51, '_', '_', 8, 34, '_', '_', 10, 4, '_']
get_hashed_value(17) -> 37 -> x mod 13 -> 11
Attempting to insert 17 at index 11
Collision at 11, which contains the value: 4
Attempting to insert 17 at index 12
Inserted 17 at index 12
1 collisions while inserting 17
Current value of hash table: [5, 25, 18, 51, '_', '_', 8, 34, '_', '_', 10, 4, 17]

Inserting 21 into [5, 25, 18, 51, '_', '_', 8, 34, '_', '_', 10, 4, 17]
get_hashed_value(21) -> 45 -> x mod 13 -> 6
Attempting to insert 21 at index 6
Collision at 6, which contains the value: 8
Attempting to insert 21 at index 7
Collision at 7, which contains the value: 34
Attempting to insert 21 at index 8
Inserted 21 at index 8
2 collisions while inserting 21
Current value of hash table: [5, 25, 18, 51, '_', '_', 8, 34, 21, '_', 10, 4, 17]
\end{verbatim}

\pagebreak

I wrote the following script to generate the trace above.

\inputminted[breaklines, breakanywhere, fontsize=\footnotesize, frame=single, fontsize=\footnotesize]{python}{problem1-hashing.py}

\pagebreak
\noindent\textbf{Problem 2:}
\medskip

Re-hashing - For this question, use the re-hashing approach to handle collisions. The re-hash algorithm is simply \texttt{get\_hashed\_value(current\_index)}, where \texttt{current\_index} is the index of the current collision.

\medskip
\noindent\textbf{Answer:}
\medskip

\begin{verbatim}
Inserting 5 into ['_', '_', '_', '_', '_', '_', '_', '_', '_', '_', '_', '_', '_']
get_hashed_value(5) -> 13 -> x mod 13 -> 0
Attempting to insert 5 at index 0
Inserted 5 at index 0
0 collisions while inserting 5
Current value of hash table: [5, '_', '_', '_', '_', '_', '_', '_', '_', '_', '_', '_', '_']

Inserting 4 into [5, '_', '_', '_', '_', '_', '_', '_', '_', '_', '_', '_', '_']
get_hashed_value(4) -> 11 -> x mod 13 -> 11
Attempting to insert 4 at index 11
Inserted 4 at index 11
0 collisions while inserting 4
Current value of hash table: [5, '_', '_', '_', '_', '_', '_', '_', '_', '_', '_', 4, '_']

Inserting 25 into [5, '_', '_', '_', '_', '_', '_', '_', '_', '_', '_', 4, '_']
get_hashed_value(25) -> 53 -> x mod 13 -> 1
Attempting to insert 25 at index 1
Inserted 25 at index 1
0 collisions while inserting 25
Current value of hash table: [5, 25, '_', '_', '_', '_', '_', '_', '_', '_', '_', 4, '_']

Inserting 8 into [5, 25, '_', '_', '_', '_', '_', '_', '_', '_', '_', 4, '_']
get_hashed_value(8) -> 19 -> x mod 13 -> 6
Attempting to insert 8 at index 6
Inserted 8 at index 6
0 collisions while inserting 8
Current value of hash table: [5, 25, '_', '_', '_', '_', 8, '_', '_', '_', '_', 4, '_']

Inserting 10 into [5, 25, '_', '_', '_', '_', 8, '_', '_', '_', '_', 4, '_']
get_hashed_value(10) -> 23 -> x mod 13 -> 10
Attempting to insert 10 at index 10
Inserted 10 at index 10
0 collisions while inserting 10
Current value of hash table: [5, 25, '_', '_', '_', '_', 8, '_', '_', '_', 10, 4, '_']

Inserting 34 into [5, 25, '_', '_', '_', '_', 8, '_', '_', '_', 10, 4, '_']
get_hashed_value(34) -> 71 -> x mod 13 -> 6
Attempting to insert 34 at index 6
Collision at 6, which contains the value: 8
Calculating new index: get_hashed_value(6) -> 15 -> x mod 13 -> 2
Attempting to insert 34 at index 2
Inserted 34 at index 2
1 collisions while inserting 34
Current value of hash table: [5, 25, 34, '_', '_', '_', 8, '_', '_', '_', 10, 4, '_']

Inserting 18 into [5, 25, 34, '_', '_', '_', 8, '_', '_', '_', 10, 4, '_']
get_hashed_value(18) -> 39 -> x mod 13 -> 0
Attempting to insert 18 at index 0
Collision at 0, which contains the value: 5
Calculating new index: get_hashed_value(0) -> 3 -> x mod 13 -> 3
Attempting to insert 18 at index 3
Inserted 18 at index 3
1 collisions while inserting 18
Current value of hash table: [5, 25, 34, 18, '_', '_', 8, '_', '_', '_', 10, 4, '_']

Inserting 51 into [5, 25, 34, 18, '_', '_', 8, '_', '_', '_', 10, 4, '_']
get_hashed_value(51) -> 105 -> x mod 13 -> 1
Attempting to insert 51 at index 1
Collision at 1, which contains the value: 25
Calculating new index: get_hashed_value(1) -> 5 -> x mod 13 -> 5
Attempting to insert 51 at index 5
Inserted 51 at index 5
1 collisions while inserting 51
Current value of hash table: [5, 25, 34, 18, '_', 51, 8, '_', '_', '_', 10, 4, '_']

Inserting 17 into [5, 25, 34, 18, '_', 51, 8, '_', '_', '_', 10, 4, '_']
get_hashed_value(17) -> 37 -> x mod 13 -> 11
Attempting to insert 17 at index 11
Collision at 11, which contains the value: 4
Calculating new index: get_hashed_value(11) -> 25 -> x mod 13 -> 12
Attempting to insert 17 at index 12
Inserted 17 at index 12
1 collisions while inserting 17
Current value of hash table: [5, 25, 34, 18, '_', 51, 8, '_', '_', '_', 10, 4, 17]

Inserting 21 into [5, 25, 34, 18, '_', 51, 8, '_', '_', '_', 10, 4, 17]
get_hashed_value(21) -> 45 -> x mod 13 -> 6
Attempting to insert 21 at index 6
Collision at 6, which contains the value: 8
Calculating new index: get_hashed_value(6) -> 15 -> x mod 13 -> 2
Attempting to insert 21 at index 2
Collision at 2, which contains the value: 34
Calculating new index: get_hashed_value(2) -> 7 -> x mod 13 -> 7
Attempting to insert 21 at index 7
Inserted 21 at index 7
2 collisions while inserting 21
Current value of hash table: [5, 25, 34, 18, '_', 51, 8, 21, '_', '_', 10, 4, 17]
\end{verbatim}

\pagebreak

As with Problem 1, I wrote a script to generate the trace above.

\inputminted[breaklines, breakanywhere, fontsize=\footnotesize, frame=single, fontsize=\footnotesize]{python}{problem2-rehashing.py}

\pagebreak

\section*{Source Code}
The full source code, including LaTeX, tooling, other artifacts, and git history is available on my Github repository for the class here:
\begin{center}
  \url{https://github.com/jllovet/jhu-en-605-202-su26-data-structures}
\end{center}

\printbibliography

\end{document}
